\scanpage{55}
Le cas $1^\circ)$ \deleted{\unsure{doit|dit}} signifie aussi que $T_\ell(E)(\beta)$ est tel que $\pi_1(X)$ y opère à travers un \uline{quotient} \uline{fini} [Mais cet énoncé ne se généralise guère via filtration de $\pi_1$, comme on voit sur $X=\operatorname{Spec}$ corps fini].

Supposons $E$ \inserted{$\in\operatorname{Ob}M^+$} \uline{semi-simple}, donc $T_\ell(E)$ semi-simple. Soit $E_i$ la composante de $E$ de type dimensionnel $i$, \deleted{\unsure{qu}} comment caractériser $T_\ell(E_i)\subset T_\ell(E)$ du pt de vue Galoisien ?
\inserted{Lorsque $X=\operatorname{Spec}$ corps fini, c'est purement \unsure{arithmétique|combinatoire} en termes de Frobenius.}
\deleted{On \unsure{pt|peut} supposer $E$ homogène de degré $i$ [la décomposition de $T_\ell(E)$ en composantes \deleted{\unsure{de degr.}} homogènes est canoniquement déterminée]. On est ramené de proche en proche \deleted{[tordant par \unsure{$\mathbb Z$} $\mathbb Z(1)$ et} [\unsure{dévissant|diminuant} par le type dimensionnel précédent] à trouver le morceau de type \inserted{dimensionnel} $i$, qd on a de type dimensionnel \deleted{$0$ ou de type dimensionnel $1$ (alors de type dimensionnel $0$ (qui sont en \unsure{qui sont nuls}}}
Le cas général s'en déduit par image inverse, si on connaît \uline{tous} les Frobenius $\in\pi_1$ \unsure{au}
% archive sheet number: 54

